\section*{Hoofdstuk \ref{ch:g2:shapes} --- Vlakke figuren en lichamen}

\begin{solution}{exo:g2:shapes:1}
$5$ zijden, $5$ hoeken: een vijfhoek. $3$ hoeken (en $3$ zijden): een
driehoek. Rond: een cirkel.
\end{solution}

\begin{solution}{exo:g2:shapes:2}
\omterm{def:g2:shapes:sort}{Veelhoeken}: vierkant, driehoek, zeshoek (rechte zijden). Geen \omterm{def:g2:shapes:sort}{veelhoek}:
de cirkel (gebogen rand).
\end{solution}

\begin{solution}{exo:g2:shapes:3}
Rechte hoeken zitten bijvoorbeeld in de hoeken van het bord, van een
boek of van een raam; een hoek die niet recht is, vind je bij het open
deksel van een doos of bij de wijzers van een klok om $2$ uur.
\end{solution}

\begin{solution}{exo:g2:shapes:4}
Vierkant: $4$ rechte hoeken. Rechthoek: $4$. Een driehoek kan hoogstens
$1$ rechte hoek hebben --- bij twee zouden de andere twee zijden elkaar
nooit meer snijden.
\end{solution}

\begin{solution}{exo:g2:shapes:5}
De kopieën op het raster maak je door de vakjes tussen de hoekpunten te
tellen; vergelijken met het voorbeeld is de controle.
\end{solution}

\begin{solution}{exo:g2:shapes:6}
Dobbelsteen: \omterm{def:g2:shapes:solids}{kubus}. Blikje soep: \omterm{def:g2:shapes:solids}{cilinder}. Voetbal: bol. Schoenendoos:
balk.
\end{solution}

\begin{solution}{exo:g2:shapes:7}
Een \omterm{def:g2:shapes:solids}{kubus} heeft $6$ \omterm{def:g2:shapes:solids}{zijvlakken}, elk een vierkant. Een balk heeft ook $6$
\omterm{def:g2:shapes:solids}{zijvlakken} (rechthoeken).
\end{solution}

\begin{solution}{exo:g2:shapes:8}
De bol rolt (hij is overal rond) en de \omterm{def:g2:shapes:solids}{cilinder} rolt (over zijn gebogen
zijkant). De \omterm{def:g2:shapes:solids}{kubus} en de balk rollen niet: zij hebben alleen vlakke
\omterm{def:g2:shapes:solids}{zijvlakken} en blijven daarop liggen.
\end{solution}

\begin{solution}{exo:g2:shapes:9}
$6$ zijden: een zeshoek, met $6$ hoeken.
\end{solution}

\begin{solution}{exo:g2:shapes:10}
De buitenrand van het huis (een vierkante muur met een driehoekig dak,
waarvan de basis op de bovenzijde van het vierkant ligt) heeft $5$ zijden
en $5$ hoeken --- het is een vijfhoek (onderzijde, twee muren, twee
dakschuintes; de bovenkant van de muur zit onder het dak verborgen).
\end{solution}
